\section{Introduction}\label{sec:introduction}

Coarse assembly relates the topology visible at bounded scales in a metric
space to the indices of operators controlled by its large-scale geometry.
Its rational injectivity asks whether a non-torsion topological class can
have zero coarse index. We prove that this can occur even when the analytic
target is the maximal Roe algebra and the space is a coarse disjoint union
of finite graphs of bounded degree.

Let $Z$ be a uniformly discrete proper metric space of bounded geometry;
thus every ball of a fixed radius contains a uniformly bounded number of
points. Write $P_r(Z)$ for the Rips complex whose simplices are the finite
subsets of diameter at most $r$, and set
\[
 \KX_q(Z)=\varinjlim_{r\to\infty}\KK_q(C_0(P_r(Z)),\C).
\]
All K-theory and K-homology in this article are complex. On the ample module
$H_Z=\ell^2(Z)\otimes\ell^2(\N)$, let $\C[Z]$ be the algebra of bounded,
locally compact, finite-propagation operators. Local compactness means that
$aT$ and $Ta$ are compact for every $a\in C_0(Z)$, and finite propagation
means that $T_{xy}=0$ whenever $d(x,y)$ exceeds some fixed bound. In
particular, the compact matrix entries of $T$ need not satisfy a uniform
rank bound. The universal C*-completion is the \emph{maximal Roe algebra}
$C^*_{\max}(Z)$. The canonical maximal coarse assembly map is
\[
 \mu_Z^{\max}:\KX_q(Z)\longrightarrow K_q(C^*_{\max}(Z)).
\]
We recall the needed norm estimates and the index construction in
Section~\ref{sec:maximal-index}.

A \emph{coarse disjoint union} of finite graphs is a metric on their disjoint
union which restricts to the graph path metric on each component and for
which only finitely many pairs of distinct components are within any fixed
distance. We take such metrics to be uniformly discrete.

\begin{theorem}\label{thm:main}
There exist finite connected graphs $(X_j)_{j\geq1}$ with a common finite
bound on vertex degree, a coarse disjoint union $X=\bigsqcup_{j\geq1}X_j$,
and an infinite-order class $\alpha\in\KX_1(X)$ such that
\[
 \mu_X^{\max}(\alpha)=0.
\]
Consequently, $\mu_X^{\max}\otimes\Q$ is not injective.
\end{theorem}

The theorem strengthens the reduced coarse-assembly result proved in the
companion paper~\cite{OpenAICoarseNovikov2026}.
If $C^*(X)$ denotes the operator-norm
closure of $\C[X]$, the canonical quotient
$C^*_{\max}(X)\to C^*(X)$ carries the maximal index to the reduced index.
Our class $\alpha$ therefore also lies in the reduced assembly kernel. The stronger
conclusion requires a proof in the universal norm: vanishing of a reduced
index alone does not imply vanishing of its maximal lift.

\subsection{Context and prior work}

Roe developed coarse cohomology and coarse index theory as tools for
large-scale geometry and index-theoretic questions on noncompact manifolds
\cite{Roe93}. Higson and Roe formulated and studied the coarse Baum--Connes
conjecture in this framework \cite{HR95}. The connection with Novikov-type
questions is that injectivity retains the topological information carried
by higher indices. The present theorem concerns the non-equivariant coarse
assembly map displayed above.

Several fundamental positive results relate assembly to geometric control.
Yu proved coarse assembly isomorphism for bounded-geometry spaces of finite
asymptotic dimension \cite[Theorem~7.1]{Yu98} and established the coarse
Baum--Connes conjecture for
bounded-geometry spaces admitting a uniform embedding into Hilbert space
\cite{Yu00}. Skandalis, Tu, and Yu developed the coarse-groupoid approach
and its relation to these embedding methods \cite{STY02}. These results
explain the importance of geometric hypotheses when seeking injectivity or
isomorphism theorems.

Earlier failures of injectivity also show why bounded geometry and the
domain of assembly matter. Yu constructed a coarse-assembly kernel class
for a space without bounded geometry \cite[Section~8, Proposition~8.1]{Yu98}.
Dranishnikov, Ferry, and Weinberger constructed integral index-kernel
examples with locally finite K-homology as domain; their kernel classes
disappear under coarsening \cite[Section~6, Remark~6.2]{DFW03}.
Here the class belongs to the actual Rips colimit for a bounded-geometry
graph union and remains nonzero after rationalization.

Finite quotients and expanding graphs also exposed limitations of general
assembly assertions. Higson, Lafforgue, and Skandalis constructed
counterexamples in group and groupoid settings and proved reduced coarse
nonsurjectivity for certain unions of congruence Cayley graphs
\cite[Section~6]{HLS02}. Nonsurjectivity is a different obstruction from
the kernel in Theorem~\ref{thm:main}. Their separate Hausdorff-groupoid
kernel construction combines an odd Bott class with a Kazhdan projection:
a nonzero full-assembly image becomes zero after reduction
\cite[Section~5]{HLS02}. This is a methodological predecessor, but its
behavior also makes it essential to distinguish the two analytic targets.

The maximal Roe algebra and its assembly map were developed in work of
Gong, Wang, and Yu \cite{GWY08} and Oyono-Oyono and Yu \cite{OOY09}; the
latter analyzed maximal assembly for specific families of expanders arising
from finite group quotients. Willett and Yu further studied higher indices
for expanding and large-girth graph families \cite{WY12}.
\v{S}pakula and Willett proved that the maximal-to-reduced comparison
induces a K-theory isomorphism for spaces admitting a coarse embedding into
Hilbert space \cite{SW13}. Chen, Wang, and Yu proved the maximal coarse
Baum--Connes conjecture under the more general hypothesis of a fibred
coarse embedding into Hilbert space \cite{CWY13}. More recent injectivity
results concern finite-group extensions whose kernel and quotient families
admit such fibred embeddings \cite{DGWY25}; maximal isomorphism results also
hold for coarse fibrations with uniformly property-A fibers, a
fibred-coarsely-embeddable base, and uniform coarse equivalences between
fibers and preimages of fixed-radius balls \cite{GWZ25}. Finite fibers alone do not
supply these uniform hypotheses. Theorem~\ref{thm:main}
shows that rational injectivity cannot hold for all bounded-geometry graph
unions without further assumptions.

Our analytic tools belong to this established maximal-index framework.
We give the norm and functoriality arguments needed here, with particular
attention to the points where the universal and operator norms must be
distinguished; see also \cite{OOY09,SW13,SS21}. Kazhdan's property~(T)
\cite{Kazhdan67,BHV08} supplies averaging in the full group C*-algebra.
The topological part uses integral line-bundle classification, the
K-theory of projective spaces, and the standard K-homology pairing
\cite{Atiyah67,Blackadar98,HigsonRoe00}.

\subsection{The construction and the proof}

The base graphs are three-dimensional discrete tori $B_j$ whose side
lengths tend to infinity. Over each $B_j$ we place a finite quotient of
$\SL_3(\Z)$, obtaining a graph $Y_j$ with Cayley edges inside the fibers.
Transport around the first two base circles is given by commuting elements
of that quotient. These elements cease to commute at a higher congruence
level. The graph $X_j$ joins two copies of $B_j$ through $Y_j$: its middle
has a long product collar, and toward either end the fiber quotients are
reduced one level at a time until they become trivial. This keeps all
vertex degrees bounded. Section~\ref{sec:geometry} gives the construction.

Write $B=\bigsqcup_jB_j$, $Y=\bigsqcup_jY_j$, and let
\[
 f:Y\longrightarrow B,\qquad i_\pm:B\longrightarrow X
\]
be the middle projection and the two end inclusions. A torus K-homology
class on each component defines a single class $\beta\in\KX_1(B)$.
The class of interest is
\[
 \alpha=(i_+)_*\beta-(i_-)_*\beta.
\]
Its persistence and its zero index come from different features of the
middle projection.

For persistence, take a line bundle $\mathcal L$ on each continuous base
torus, with first Chern class the product of the first two circle classes.
Its reduced K-theory class, multiplied by the third circle's odd class,
pairs to one with the chosen K-homology class. At every fixed Rips scale,
the lifted commutator calculation shows that the pullback of $\mathcal L$
to the middle admits a $2^j$-th root for all sufficiently large $j$.
A uniform dimension bound for the Rips complexes converts these roots
into divisibility by a growing power of two in integral K-theory.
Mayer--Vietoris then prevents any fixed nonzero multiple of the endpoint
difference from vanishing. Section~\ref{sec:persistence} makes this
argument before taking either the scale colimit or rational coefficients.

For the index, projection from the middle induces a surjection
\[
 f_*:K_1(C^*_{\max}(Y))\longrightarrow K_1(C^*_{\max}(B)).
\]
To construct a splitting, we transport each base matrix entry horizontally
between the fibers and then apply the Kazhdan projection onto the constant
vectors in each fiber. The transport lift need not preserve products;
averaging makes it multiplicative, and projecting the resulting
homomorphism back to the base induces the identity on K-theory.
Section~\ref{sec:transfer} proves these identities in the maximal completion
using two controls from Section~\ref{sec:maximal-index}: agreement of
maximal and operator norms on finite vertex supports, and uniform norm
estimates at fixed propagation.

Finally, Section~\ref{sec:sliding} shows that the two composites
$(i_+)_*f_*$ and $(i_-)_*f_*$ agree on analytic $K_1$. The end maps are
not uniformly close: their separation grows with $j$. Instead, we place
orthogonal copies of an operator at every height and rotate these copies
simultaneously by one vertical step. After cancelling the common interior
copies, the endpoint maps agree modulo compact operators, and injectivity
of $K_1(C^*_{\max}(X))\to K_1(C^*_{\max}(X)/\K)$ removes that quotient.
Surjectivity of $f_*$ and naturality of assembly then give the zero index.
